% Pursuing Stacks (cotes 134-2 to 134-8): archive cote 134-2, site page 69 (PDF page 70).
% appendix (letters to Breen). Draft compared: ps1.tex:1563-1570.
% AI-transcribed by Claude Opus 5.5 (claude-opus-5-5, adaptive thinking, effort high), Message Batches API, batch ps1 (msgbatch_01Rv5hCaNTtUt4JHx1yqoHPL), 2026-11-12; not yet checked by eye.
\scanpage{134-2/69}
\mspage{-54-}

feels like saying: which commute with finite $\varprojlim$ and arbitrary $\varinjlim$); but
this raises the disquieting vision of n-limits in n-categories. (N.B. The
case $n = 2$ begins to become famliliar to us ...). It is prudent in all
of this to suppose that $X$ is ``locally homotopically trivial'', which
ensures the pro-simplicial set which Artin-Mazur associate to it (with
the help of nerves of hyper-coverings) is essentially constant in the
ordinary homotopy category - thus $X$ defines a homotopy type in the usual
sense. This is surely \textit{not} the case for the étale topos of a scheme. In
such case, the fundamental n-groupoid should be conceived as a \uline{\textit{pro}}-n-\textit{groupoid}
(nothing surprising in that, in view of the familiar theory of $\pi_1$), and
$E_n$ as an (ind)-n-category (the ind-structure will correspond to the
exigencies of local triviality for a variable n-stack, relative to coverings
more and more fine on $X$) .

\rule{2cm}{0.4pt}\$\rule{2cm}{0.4pt}

\margin{(App.\ 16)} I nevertheless understand your instinctive resistance to conceive
this extreme stripping of a beautiful topos $X$ , to the point of retaining
only the meagre homotopy type. Even more\hand{,} I am persuaded that going to the
root of this instinctive resistance, one arrives at a generalisation and
deepening of the notion of ``homotopy type'', and to bring new grist to the
mill of the development of a good homotopical yoga. Here is what I have
in mind.

Let us speak first of sheaves (of sets, or of modules, etc.) instead
of stacks, for simplicity, and place ourselves in the étale topos of a
scheme. The locally constant sheaves - modulo a supplementary condition
of finiteness which is sufficiently anodyne - form the easiest of the
\textit{constructible} sheaves, for the definition of which they serve as models.
Supposing $X$ coherent (= quasi-compact and quasi-separated), then the
general constructible sheaves are those for which there exist a finite
partition $X = \bigcup_{i\in I} X_i$ of $X$ into ``cells'' or ``strata'' $X_i$ , each locally

\margin{69}
